% Generated draft: requires mathematical and source review.
\documentclass[11pt]{article}
\usepackage[margin=25mm]{geometry}
\usepackage{fontspec}
\setmainfont{Latin Modern Roman}
\usepackage{amsmath,amssymb,mathtools,unicode-math}
\setmathfont{Latin Modern Math}
\usepackage{enumitem,needspace,xurl,hyperref}
\setlength{\parindent}{0pt}
\setlength{\parskip}{5pt}
\setlength{\emergencystretch}{3em}
\newcommand{\N}{\mathbb{N}}
\newcommand{\Z}{\mathbb{Z}}
\newcommand{\Q}{\mathbb{Q}}
\newcommand{\R}{\mathbb{R}}
\newcommand{\C}{\mathbb{C}}
\newcommand{\F}{\mathbb{F}}
\newcommand{\A}{\mathbb{A}}
\newcommand{\PP}{\mathbb{P}}
\newcommand{\E}{\mathbb{E}}
\newcommand{\eps}{\varepsilon}
\newcommand{\dd}{\,\mathrm{d}}
\begin{document}
\section*{UnsolvedMath ID: 30000141}
\textbf{Source problem code:} catalogue entry
\subsection{Definitions and mathematical statement}
Use the usual Euclidean constant-mean-curvature (CMC) surface setting. A coplanar \(k\)-unduloid is understood in the standard complete Alexandrov-embedded genus-zero, \(k\)-ended convention, with unduloid-type ends and a common reflection plane; the exact source class needs verification. The neck parameters \(n_i\) are not defined in the excerpt, so their scale must be fixed consistently with its displayed equalities.

For a two-sided CMC immersion \(F:M\to\R^3\), let \(A\) be the second fundamental form and use Jacobi operator
\[
 J_M=\Delta_M+|A|^2.
\]
A scalar normal Jacobi field satisfies \(J_Mu=0\). Define
\[
 \mathcal J_{L^2}(M)=\{u\in C^\infty(M):J_Mu=0,\ \int_M|u|^2\,\dd A_M<\infty\}.
\]
The sign convention for the Laplacian is the usual geometric one, with the equivalent kernel unaffected by a simultaneous operator sign change.

The source has four distinct questions:
\begin{enumerate}
\item Are the stated upper bounds on \(\dim\mathcal J_{L^2}(M)\) sharp for coplanar \(k\)-unduloids? The bound's formula, dependence on neck data, and whether sharpness is pointwise or over a family are not supplied.
\item For a triunduloid, does a nonzero \(L^2\) Jacobi field exist in either exceptional neck configuration
\[
 n_1+n_2+n_3=2\pi,
 \qquad n_i+n_j=n_\ell\quad\text{with }\{i,j,\ell\}=\{1,2,3\}?
\]
The source's normalization and the existence/admissibility of surfaces at these equalities must be specified; no positive kernel dimension is inferred merely from equality.
\item Can the methods extend to other CMC surfaces to bound neck sizes or the dimensions of Jacobi-field spaces? The methods, target surface class, and field growth conditions are not given.
\item Independently, is every tempered Jacobi field on a CMC surface integrable? Here integrability would mean realization as the normal derivative of a differentiable family of CMC immersions, with the fixed-mean-curvature and admissible surface/deformation class declared. ``Tempered'' requires an end-growth or weighted-space definition; it is not silently identified with \(L^2\), bounded, or one fixed polynomial growth bound.
\end{enumerate}
The last question is not reduced to the first three or restricted to a neck equality by the source. Modding out ambient motions, fixing topology/end data, and varying the mean curvature can change the formulation and require explicit conventions.

The supplied literature triage dated 21 August 2026 reports broad later nondegeneracy results for coplanar CMC surfaces, but says the report's combined sharpness, exceptional-case, extension, and general integrability questions are not all resolved. It does not specify exactly which hypotheses and subquestions the cited theorem settles. No current existence or integrability claim for an exceptional case is inferred from that broad assessment.
\subsection{Short English statement}
The draft preserves four separate targets: sharpness of missing dimension bounds, \(L^2\)-field existence at the two stated triunduloid equalities, extensions to other CMC surfaces, and integrability of tempered fields. Neck normalization, growth conditions, and deformation classes remain necessary source details.
\subsection{Sources}
\begin{itemize}
\item Local UnsolvedMath registry, permanent ID \(30000141\), exact source code \texttt{catalogue entry}, Oberwolfach Reports Open Problems. Statement status \texttt{corrected\_verified}; the supplied verification reports direct source checking and faithful separation of the four questions.
\item ``Geometrie (2005),'' Oberwolfach Reports, pp. 2491--2538, DOI \texttt{10.4171/owr/2004/47}: \url{https://doi.org/10.4171/owr/2004/47}. Proposed year and DOI refer to 2004, while the citation label says 2005.
\item Supplied ``coplanar k-unduloids paper'': \url{https://doi.org/10.1093/imrn/rnp058}. The other supplied source is the original report; exact upper bounds, neck definitions, and nondegeneracy theorem hypotheses are not reproduced.
\item Literature checked 21 August 2026; supplied status \texttt{partially\_solved}. Registry created 25 August 2026 and updated 16 September 2026. This dated curation triage requires independent verification before citation.
\item Registry curation attribution: Ulam AI / UnsolvedMath Contributors, CC BY 4.0. Underlying report and literature retain their supplied provenance; this statement-only draft requires human review.
\end{itemize}
\end{document}
